\section{Reciprocal Green functions and a squared-distance kernel}\label{sec:reciprocal}

In this section we prove the reciprocal-kernel theorem for a general disk potential. The kernels $K,N$ are those of \eqref{eq:K}--\eqref{eq:N}; their construction and estimates require only the measure hypotheses restated in the theorem below.

A real symmetric kernel $D$ with zero diagonal is \emph{conditionally negative semidefinite} if, for every finite family of points and real coefficients with $\sum_i a_i=0$,
\[
\sum_{i,j}a_i a_jD(s_i,s_j)\le0.
\]

\begin{theorem}\label{thm:negative-kernel}
Let $\dd q=a\dd x$ on $\D$, where $a\in L^\infty(\D)$ is nonnegative, and suppose there is a constant $0<d<1$ such that
\[
\int_\D|\psi|^2\dd q\le d\int_\D|\grad\psi|^2
\qquad(\psi\in H^1_0(\D)).
\]
For the associated kernels $K,N$ and every $p\in\D$, the kernel
\begin{equation}\label{eq:D}
D_p(s,t)=\frac{K(p,s)K(p,t)}{N(s,t)}\quad(s\ne t),
\qquad D_p(s,s)=0,
\end{equation}
is conditionally negative semidefinite on $\Sone$.
\end{theorem}

\begin{remark}\label{rem:zero-potential}
For $q=0$, this kernel is a rescaled squared chordal distance. Indeed, with the disk automorphism $\varphi_p(z)=(z-p)/(1-\overline pz)$,
\[
D_p(s,t)=\pi P_s(p)P_t(p)|s-t|^2
=\frac1{4\pi}|\varphi_p(s)-\varphi_p(t)|^2.
\]
\Cref{thm:negative-kernel} supplies a Hilbert-space squared-distance interpretation after a nonnegative subcritical potential is added.
\end{remark}

We prove the theorem through reciprocal matrices. The reciprocal inertia condition comes from the McCullough--Quiggin theory \cite{McCullough1994,Quiggin1994}. For finite-valued positive-definite kernels whose entries never vanish, it characterizes the complete Nevanlinna--Pick property \cite[Corollary~1.12]{AglerMcCarthy2000}. Quiggin established the reciprocal condition for a class of weighted one-dimensional Sobolev kernels, represented by Green functions of regular Sturm--Liouville problems \cite[Section~2.2, Corollary~2.2.2 and Theorem~2.2.3]{Quiggin1994}. Here this condition serves as an algebraic model; all the matrix and analytic arguments needed for the singular Green kernel will be given below.

Say that a real symmetric matrix has property \textup{(P)} if it has at most one positive eigenvalue, counted with multiplicity. This property is preserved by positive diagonal congruence, subtraction of a positive semidefinite matrix, passage to principal submatrices, and limits of matrices of fixed size. Reciprocals below are always \emph{entrywise}, never matrix inverses.

The fixed interior point $p$ connects \textup{(P)} to the desired conditional negativity. Adjoining it to boundary points will give a reciprocal matrix with zero diagonal, boundary block $B_{ij}=1/N(s_i,s_j)$ for $i\ne j$, and border $\ell_i=1/K(p,s_i)$. Property \textup{(P)} for this bordered matrix forces $x^TBx\le0$ on $\ell^Tx=0$; the substitution $x_i=K(p,s_i)a_i$ turns this into the required inequality for $\sum_i a_i=0$. We prove this last algebraic step at the end of the section.

We first construct bounded regularizations of $G$ whose reciprocals have \textup{(P)}, and prove that updates along a column of the current kernel matrix preserve this condition. The resolvent construction then uses four limits, in order: refine a positive quadrature at fixed $\epsilon>0$ and compactly supported $0\le\rho\le q$; let $\epsilon\downarrow0$; send the moving evaluation points radially to the circle while fixing $p$; and finally increase $\rho$ to $q$. Compact support keeps the Poisson sections bounded during the boundary step.

\subsection{A bounded regularization of the Green kernel}

Direct quadrature of $G$ would encounter the infinite value $G(z,z)$ at every integration node. We therefore need a bounded replacement with finite positive diagonal that retains the reciprocal condition \textup{(P)}.

The scalar integral used below is an equivalent form of the logarithmic-mean representations in \cite{QiZhangLi2014}. We derive it and then establish the positive-semidefinite disk-kernel property that we need.

\begin{lemma}\label{lem:regularization}
There are bounded continuous strictly positive kernels $C_\epsilon$ on $\D\times\D$, $0<\epsilon\le1$, such that $C_\epsilon\uparrow G$ as $\epsilon\downarrow0$, including on the diagonal, and every finite matrix $(1/C_\epsilon(x_i,x_j))_{i,j}$ has \textup{(P)}.
\end{lemma}
\begin{proof}
Put
\[
h(x,y)=\frac{|x-y|^2}{(1-|x|^2)(1-|y|^2)},\qquad
L(h)=\frac1{\log(1+1/h)}\ (h>0),\quad L(0)=0,
\]
and $f(h)=h+\tfrac12-L(h)$. After multiplying the $x$- and $y$-entries by $1-|x|^2$ and $1-|y|^2$, respectively, the kernel $h+\tfrac12$ becomes
\begin{equation}\label{eq:lorentz}
\tfrac12(1+|x|^2)(1+|y|^2)-2x\cdot y.
\end{equation}
This is a positive rank-one kernel minus a positive semidefinite kernel, so $h+\tfrac12$ has \textup{(P)} on finite sets.

Since $1/G=4\pi L(h)$, the identity $L=(h+\tfrac12)-f$ will be useful once we prove that $f(h(x,y))$ is positive semidefinite. For $h>0$,
\begin{equation}\label{eq:log-mean}
L(h)=\int_0^1h^{1-\alpha}(h+1)^\alpha\dd\alpha.
\end{equation}
For $0<\alpha<1$, let $c_\alpha$ normalize $v^{\alpha-1}(1+v)^{-1}\dd v$ to a probability measure on $(0,\infty)$, and set $r=(1+v)^{-1}$. Integration by parts gives
\[
\mathbb E_\alpha(1-r)
=c_\alpha\int_0^\infty\frac{v^\alpha}{(1+v)^2}\dd v
=\alpha.
\]
For $h>0$, scaling $v$ in the normalized integral yields
\begin{align}
h^{1-\alpha}(h+1)^\alpha
&=c_\alpha\int_0^\infty
\frac{h(h+1)}{(h+1)+vh}v^{\alpha-1}\dd v\notag\\
&=\mathbb E_\alpha\frac{h(h+1)}{h+r}
=\mathbb E_\alpha\left[h+(1-r)-\frac{r(1-r)}{h+r}\right].\label{eq:mixture-algebra}
\end{align}
The last equality remains true at $h=0$. Integrating in $\alpha$ therefore gives
\begin{equation}\label{eq:f-mixture}
f(h)=\int_0^1\mathbb E_\alpha\frac{r(1-r)}{h+r}\dd\alpha,
\qquad 0\le f(h)\le f(0)=\tfrac12.
\end{equation}

Now
\[
k(x,y)=\frac1{1+h(x,y)}
=\frac{(1-|x|^2)(1-|y|^2)}{|1-x\overline y|^2}
\]
is positive semidefinite. Indeed $(1-x\overline y)^{-1}$ is a Gram kernel by its power series; multiplying it by its complex conjugate and the positive diagonal factors preserves positive semidefiniteness. For fixed $0<r<1$,
\begin{equation}\label{eq:schur-geometric}
\frac1{h+r}=\sum_{n\ge0}(1-r)^n k^{n+1},
\end{equation}
where powers are entrywise. This series converges even when $h=0$. Each power is positive semidefinite, as follows by taking tensor products of Gram vectors. Thus $1/(h+r)$ is positive semidefinite, and the positive mixture \eqref{eq:f-mixture} proves the claim. The mixture has finite entries by the displayed bound, including on the diagonal.

Define
\begin{equation}\label{eq:C-epsilon}
\frac1{C_\epsilon(x,y)}
=4\pi\bigl[h+\tfrac12-(1-\epsilon)f(h)\bigr]
=4\pi\bigl[L(h)+\epsilon f(h)\bigr].
\end{equation}
Its reciprocal matrix has \textup{(P)} by \eqref{eq:lorentz} and the positive semidefiniteness just proved. Its denominator is at least $2\pi\epsilon$, because
\[
L+\epsilon f=\epsilon(h+\tfrac12)+(1-\epsilon)L\ge\epsilon/2.
\]
It is positive and continuous, so $C_\epsilon$ is bounded, continuous, and strictly positive. Since $f\ge0$, these kernels increase to $G$ by \eqref{eq:green}, and $0<C_\epsilon\le G$. On the diagonal $C_\epsilon(x,x)=1/(2\pi\epsilon)$, as required.
\end{proof}

\subsection{A finite resolvent update}

\begin{lemma}\label{lem:update}
Let $M$ be a real symmetric finite matrix with strictly positive entries such that its reciprocal has \textup{(P)}. For any index $z$ and any $c\ge0$, the matrix
\[
\widetilde M_{ij}=M_{ij}+cM_{iz}M_{zj}
\]
also has a reciprocal with \textup{(P)}.
\end{lemma}
\begin{proof}
Positive diagonal congruence transforms the reciprocal of $M$ into
\[
H_{ij}=\frac{M_{iz}M_{zj}}{M_{ij}}.
\]
Writing $a=M_{zz}>0$, its $z$-row and column are identically $a$. Put $z$ last, and let $H_0$ denote the remaining principal block. Subtracting the last row from every other row, and doing the same to the columns, gives the congruence
\[
H\ \sim\
\begin{pmatrix}H_0-a\one\one^T&0\\0&a\end{pmatrix}.
\]
Because $a>0$ and $H$ has \textup{(P)}, the first block is nonpositive. Hence
\[
U=a\one\one^T-H\PSD.
\]
Since $H_{ii}>0$, one has $0\le U_{ii}<a$, and positive-semidefinite Cauchy--Schwarz gives $|U_{ij}|<a$.

Using the same diagonal normalization, the new reciprocal has entries $H_{ij}/(1+cH_{ij})$. The difference from the constant rank-one kernel is
\begin{align}
\frac a{1+ca}-\frac{H_{ij}}{1+cH_{ij}}
&=\frac{U_{ij}}{(1+ca)(1+ca-cU_{ij})}\notag\\
&=\sum_{n\ge0}\frac{c^nU_{ij}^{n+1}}{(1+ca)^{n+2}}.\label{eq:update-series}
\end{align}
The series converges absolutely because $|cU_{ij}|<1+ca$. Its coefficients are nonnegative and every entrywise power of $U$ is positive semidefinite, even if $U$ has negative entries. Thus the normalized reciprocal of $\widetilde M$ is a positive rank-one kernel minus a positive semidefinite matrix, proving \textup{(P)}. For $c=0$ the identity reduces directly to $U$.
\end{proof}

\subsection{The interior resolvent kernel}

For $0\le\rho\le q$ supported on a compact subset of $\D$, let $T_\rho$ be the Green operator on $L^2(\rho)$ and define
\begin{equation}\label{eq:M-rho}
M^\rho(x,y)=G(x,y)
+\ip{G(x,\cdot)}{(I-T_\rho)^{-1}G(y,\cdot)}_\rho.
\end{equation}
It is finite for distinct interior points, since its Green sections belong to $L^2(\rho)$, and has infinite diagonal.

\begin{proposition}\label{prop:interior-reciprocal}
For every finite set of interior points, the reciprocal matrix of $M^\rho$, with diagonal zero, has \textup{(P)}.
\end{proposition}
\begin{proof}
First fix $\epsilon>0$ and replace $G$ throughout \eqref{eq:M-rho} by $C=C_\epsilon$. Its integral operator $T_C$ on $L^2(\rho)$ satisfies
\[
|T_Cf|\le T_\rho|f|,\qquad \norm{T_C}\le d<1.
\]
The absolute pointwise domination proves the norm bound without requiring a separate positivity assertion about its quadratic form.

Partition a compact integration set into finitely many measurable cells with uniformly vanishing diameters. For each cell of positive $\rho$-mass, choose a representative, and let $\pi$ send that cell to the representative. Use its mass as the quadrature weight. Uniform continuity gives uniform convergence of
\[
C(\pi(\cdot),\pi(\cdot))\longrightarrow C
\]
on the integration set up to null sets. The associated operators $T_\pi$ on $L^2(\rho)$ converge to $T_C$ in operator norm; for example their norm difference is at most $\rho(\D)$ times the uniform kernel error. Likewise $C(x,\pi(\cdot))\to C(x,\cdot)$ in $L^2(\rho)$ for each fixed evaluation point $x$. In particular, sufficiently fine quadratures have $\norm{T_\pi}<1$, and
\begin{equation}\label{eq:quadrature-resolvent}
C(x,y)+\ip{C(x,\pi(\cdot))}{(I-T_\pi)^{-1}C(y,\pi(\cdot))}_\rho
\end{equation}
converges to the regularized continuum resolvent evaluation. This follows from convergence of the source vectors and of the inverses in operator norm.

For clarity, if the nodes and weights are $z_\alpha,w_\alpha$, set
\[
\mathsf T_{\alpha\beta}
=\sqrt{w_\alpha}\,C(z_\alpha,z_\beta)\sqrt{w_\beta},
\qquad
(\xi_x)_\alpha=\sqrt{w_\alpha}\,C(x,z_\alpha).
\]
The normalized cell indicators form an orthonormal basis for the step functions. On this subspace $T_\pi$ is represented by $\mathsf T$, and it vanishes on the orthogonal complement. Thus $\norm{\mathsf T}<1$, and the correction in \eqref{eq:quadrature-resolvent} is $\xi_x^T(I-\mathsf T)^{-1}\xi_y$. Evaluation points enter only these finite vectors, including when they coincide with nodes.

We next prove \textup{(P)} for each sufficiently fine quadrature. On the union of the evaluation points and quadrature nodes, start with $M_{ij}=C(x_i,x_j)$. By \Cref{lem:regularization}, its reciprocal has \textup{(P)}. Expanding the finite resolvent sums all chains, multiplying kernel entries along each chain and the quadrature weight at every intermediate occurrence. These nonnegative sums are finite by the preceding matrix formula. Restricting the allowed intermediate nodes gives a subseries and hence remains finite.

Suppose some nodes have been allowed and their chain sum is $M$. Admit one more node $z$ of weight $w>0$. Chains with exactly $k\ge1$ internal occurrences of $z$ have total weight
\[
w^kM_{iz}M_{zz}^{k-1}M_{zj}.
\]
All factors are positive. Finiteness of the full series forces $wM_{zz}<1$, and summing over $k$ updates the matrix to
\begin{equation}\label{eq:elimination}
M_{ij}+\frac{w}{1-wM_{zz}}M_{iz}M_{zj}.
\end{equation}
\Cref{lem:update} therefore propagates \textup{(P)} one node at a time. This counting remains valid when an endpoint is a quadrature node, since only internal occurrences receive weights. Coincident node locations may be merged by adding their weights. Thus the discrete kernel \eqref{eq:quadrature-resolvent} has the required reciprocal property. The quadrature limit gives it for the regularized continuum kernel. If $\rho=0$, this conclusion follows directly from \Cref{lem:regularization}.

Finally let $\epsilon\downarrow0$. Every chain integral increases to its counterpart with $G$ by monotone convergence, as does the sum of these nonnegative integrals. Off the diagonal, its limit is \eqref{eq:M-rho}, and the correction is finite because
\[
\big|\ip{G(x,\cdot)}{(I-T_\rho)^{-1}G(y,\cdot)}_\rho\big|
\le\frac{\norm{G(x,\cdot)}_\rho\norm{G(y,\cdot)}_\rho}{1-d}.
\]
The direct diagonal term tends to infinity, so its reciprocal tends to zero. Finite entrywise limits preserve \textup{(P)}, completing the proof.
\end{proof}

\subsection{The boundary limit}

For compactly supported $\rho$, write $K^\rho,R^\rho,N^\rho$ for the kernels in \eqref{eq:K}--\eqref{eq:N} with $q$ replaced by $\rho$. The boundary limit will be taken with this compact support fixed, before increasing $\rho$ to $q$.

\begin{lemma}\label{lem:bordered}
For distinct $s_1,\ldots,s_m\in\Sone$ and $p\in\D$, the matrix
\begin{equation}\label{eq:bordered}
\begin{pmatrix}B&\ell\\\ell^T&0\end{pmatrix},
\qquad
B_{ij}=\begin{cases}1/N(s_i,s_j),&i\ne j,\\0,&i=j,\end{cases}
\qquad \ell_i=1/K(p,s_i),
\end{equation}
has \textup{(P)}.
\end{lemma}
\begin{proof}
Fix compactly supported $0\le\rho\le q$. Apply \Cref{prop:interior-reciprocal} to the points $r s_1,\ldots,r s_m,p$, with $r\uparrow1$ and $r>|p|$. The explicit Green formula gives
\begin{align}
\frac{G(rs,x)}{1-r}&\longrightarrow P_s(x),\label{eq:radial-one}\\
\frac{G(rs,rt)}{(1-r)^2}&\longrightarrow N_0(s,t)\qquad(s\ne t).\label{eq:radial-two}
\end{align}
The first convergence is uniform on compact interior sets. For the second, one may use
\[
G(rs,rt)=\frac1{4\pi}\log\left(1+\frac{(1-r^2)^2}{r^2|s-t|^2}\right).
\]
The bounded resolvent in \eqref{eq:M-rho} and the uniform convergence on $\supp\rho$ therefore imply
\[
\frac{M^\rho(rs_i,rs_j)}{(1-r)^2}\longrightarrow N^\rho(s_i,s_j)
\quad(i\ne j),\qquad
\frac{M^\rho(p,rs_i)}{1-r}\longrightarrow K^\rho(p,s_i).
\]
Here the first resolvent correction tends to
$\ip{P_{s_i}}{(I-T_\rho)^{-1}P_{s_j}}_\rho=R^\rho(s_i,s_j)$.
The second tends to
$P_{s_i}(p)+\ip{G(p,\cdot)}{(I-T_\rho)^{-1}P_{s_i}}_\rho=K^\rho(p,s_i)$.
These pairings are legitimate on the compact support, where the Poisson kernels are bounded.

Scale the reciprocal matrix by positive diagonal congruence with factors $1-r$ at the moving points and factor $1$ at $p$. Its limit is \eqref{eq:bordered} with $K^\rho,N^\rho$, so it has \textup{(P)}. The zero diagonal stays zero throughout.

Now choose increasing disk cutoffs $\dd\rho_n=\one_{\{|x|\le r_n\}}\dd q$, with $r_n\uparrow1$. Every term of $K^{\rho_n}$ and $R^{\rho_n}$ is a nonnegative chain integral. Monotone convergence in each finite product integral and then in the series gives
\[
K^{\rho_n}(p,s)\uparrow K(p,s),\qquad
N^{\rho_n}(s,t)\uparrow N(s,t)\quad(s\ne t).
\]
These limits are finite by \Cref{lem:kernel-estimates}. Taking the reciprocals and the finite matrix limit proves \eqref{eq:bordered}. No $L^2(q)$ boundary limit of an individual Poisson kernel has been used.
\end{proof}

\begin{proof}[Proof of \Cref{thm:negative-kernel}]
Use the matrix of \Cref{lem:bordered}. We first show that $x^TBx\le0$ whenever $\ell^Tx=0$. If instead $x^TBx>0$, choose $y$ with $\ell^Ty\ne0$ and replace it by
\[
y-\frac{x^TBy}{x^TBx}x.
\]
Then $x^TBy=0$ and still $\ell^Ty\ne0$. In the bordered form, $(x,0)$ and $(y,c)$ are orthogonal, and $c$ can be chosen so that
\[
(y,c)^T\begin{pmatrix}B&\ell\\\ell^T&0\end{pmatrix}(y,c)
=y^TBy+2c\ell^Ty>0.
\]
They would span a two-dimensional positive subspace, contradicting \textup{(P)}.

Finally, for coefficients with $\sum_i a_i=0$, put $x_i=K(p,s_i)a_i$. Then $\ell^Tx=0$ and
\[
\sum_{i,j}a_i a_jD_p(s_i,s_j)=x^TBx\le0.
\]
This proves the assertion on distinct points; repeated points are handled by combining their coefficients.
\end{proof}
